% SGA 7 I，PDF 3—33；原目录 PDF 9 由中文构建重建。
\SourcePage{3}
\begin{center}
{\large 数学讲义丛书（Lecture Notes in Mathematics）}\\
非正式报告与研讨班文集\\
主编：A. Dold（Heidelberg）、B. Eckmann（Zürich）\\[1ex]
第 288 卷\\[1ex]
{\Large 代数几何中的单值群}\\[1ex]
Springer-Verlag\\
Berlin · Heidelberg · New York，1972
\end{center}
\SourcePage{4}
\noindent 美国数学学会学科分类（1970）：主要分类：14D05、14D15、14K99；次要分类：13D10、14E20、14F20。
\noindent ISBN 3-540-05987-3，Springer-Verlag Berlin–Heidelberg–New York。\\
ISBN 0-387-05987-3，Springer-Verlag New York–Heidelberg–Berlin。
本书受版权保护。无论涉及全部材料还是其中一部分，一切权利均予保留，特别包括翻译、重印、插图再利用、广播、以影印机或类似方式复制，以及存入数据库的权利。依德国著作权法第 54 条，非供私人使用而制作复制品时，须向出版社付费，费用数额由与出版社的协议确定。
\noindent © Springer-Verlag Berlin · Heidelberg，1972。美国国会图书馆编目卡号 72-88730。德国印制。\\
胶版印刷：Julius Beltz，Hemsbach/Bergstr。
\SourcePage{5}
\begin{center}
{\large Bois-Marie 代数几何讨论班}\\
1967—1969\\[1ex]
{\Large 代数几何中的单值群}\\
（SGA 7 I）\\[1ex]
主持：A. Grothendieck\\
协作：M. Raynaud、D. S. Rim
\end{center}
\SGAchapter{}{引言}{Introduction}
\SourcePage{6}
设 $S$ 为\Term{scheme}{概形}，$f:X\to S$ 为\Term{scheme}{概形}\Term{morphism}{态射}。若 $f$ \Term{proper}{固有}且\Term{smooth}{光滑}，素数 $\ell$ 在 $S$ 上可逆，则 $f$ 的\Term{geometric-fibre}{几何纤维}的 $\ell$ 进\Term{cohomology}{上同调}群 $H^i(X_{\bar x},\mathbf Z_\ell)$ 构成 $S$ 上的一个 $\ell$ 进\Term{local-system}{局部系统}。若仅假设 $f$ \Term{proper}{固有}，则这些群是 $S$ 上的一个 $\ell$ 进\Term{sheaf}{层} $R^if_*\mathbf Z_\ell$ 的\Term{fibre}{纤维}。\Term{vanishing-cycles}{消失圈}理论把此\Term{sheaf}{层}在 $S$ 上的\Term{ramification}{分歧}与 $f$ 的奇点联系起来。
我们只考虑 $S$ 为 \Term{henselian-trait}{Hensel 离散赋值环谱}（trait hensélien，即 Hensel \Term{discrete-valuation-ring}{离散赋值环}的谱）的情形。这在实际中是主要情形。关于本理论的启发式描述，参见第 I 讲 2.1。设 $S=\operatorname{Spec}(V)$，$k(\bar\eta)$ 为 $V$ 的\Term{fraction-field}{分式域} $k(\eta)$ 的一个\Term{algebraic-closure}{代数闭包}，$k(\bar s)$ 为\Term{residue-field}{剩余域} $k(s)$ 的相应\Term{algebraic-closure}{代数闭包}（$k(\bar s)$ 是 $V$ 在 $k(\bar\eta)$ 中的\Term{normalization}{正规化} $\bar V$ 的\Term{residue-field}{剩余域}）。特别地，若 $X$ 在 $S$ 上\Term{proper}{固有}、\Term{flat}{平坦}，\Term{relative-dimension}{相对维数}为 $n$，且 $f$ 只有一个非\Term{smooth}{光滑}点 $x\in X_s$，则可定义一些 $\operatorname{Gal}(\bar\eta/\eta)$ 模 $\phi^i$；它们完全取决于 $x$ 附近的局部情形，且在 $i\notin[0,n]$ 时为零（第 I 讲 4.2）。还可构造 $\operatorname{Gal}(\bar\eta/\eta)$ 模的\Term{long-exact-sequence}{长正合列}
\[
\cdots\longrightarrow H^i(X_{\bar s},\mathbf Z_\ell)
\xrightarrow{\mathrm{sp}}H^i(X_{\bar\eta},\mathbf Z_\ell)
\longrightarrow\phi^i\longrightarrow H^{i+1}(X_{\bar s},\mathbf Z_\ell)
\longrightarrow\cdots .
\]
（这里把 $\operatorname{Gal}(\bar s/s)$ 模 $H^i(X_{\bar s},\mathbf Z_\ell)$ 看作惯性\Term{action}{作用}平凡的 $\operatorname{Gal}(\bar\eta/\eta)$ 模；$\mathrm{sp}$ 为\Term{specialization}{特化}\Term{arrow}{箭头}。）我们还给出保证低次 $\phi^i$ 为零的判据；目前这些判据只在\Term{characteristic}{特征} $0$ 时成立（例如，若 $X$ 还是\Term{local-scheme}{局部}\Term{complete-intersection}{完全交}，则 $i\ne n$ 时 $\phi^i=0$，见第 I 讲 4.5）。

\SourcePage{7}
\Term{monodromy-theorem}{单值性定理}断言，$\operatorname{Gal}(\bar\eta/\eta)$ 的\Term{inertia-subgroup}{惯性子群} $I$ 的\Term{action}{作用}是\Term{quasi-unipotent}{拟幺幂}的：对 $T\in I$，存在整数 $N>0$、$M>0$，使 $H^i(X_{\bar\eta},\mathbf Z_\ell)$ 的\Term{endomorphism}{自同态} $(T^M-1)^N$ 为零。本书给出此\Term{theorem}{定理}的两个证明。第一个证明（第 I 讲 1.2）具有算术性质，只要所考虑的\Term{cohomology}{上同调}群具有适当的\Term{finiteness}{有限性}性质，它就适用。关键在于：若 $k(s)$ 在其\Term{prime-subfield}{素子域}上\Term{finite-generated}{有限生成}，则对 $\operatorname{Gal}(\bar\eta/\eta)$ 的任意 $\ell$ 进表示，$I$ 的\Term{action}{作用}都是\Term{quasi-unipotent}{拟幺幂}的。第二个证明更具几何性质，它需要\Term{purity}{纯性}和\Term{singularity-resolution}{奇点解消}；目前仅在\Term{characteristic}{特征} $0$ 时或对 $H^1$ \Term{effective}{有效}。它给出关于\Term{nilpotence-index}{幂零指数} $N$ 的宝贵信息（对 $H^i$，有 $N\le i+1$）。以后将看到，在 $\mathbf C$ 上，它很适于同\Term{transcendental}{超越}理论作比较。

把这些结果用于\Term{abelian-variety}{阿贝尔簇}的 $H^1$，即其 \Term{tate-module}{Tate 模}，就能研究这些簇的模 $p$ \Term{reduced}{约化}。这样得到\Term{abelian-variety}{阿贝尔簇}\Term{stable-reduction}{稳定约化}\Term{theorem}{定理}的两个证明。此\Term{theorem}{定理}断言：经\Term{ramification}{分歧}\Term{group-extension}{扩张}后，\Term{neron-model}{Néron 模型}的\Term{connected}{连通}\Term{special-fibre}{特殊纤维}是\Term{abelian-variety}{阿贝尔簇}被\Term{torus}{环面}的\Term{group-extension}{扩张}。第一个证明（第 I 讲第 6 节）沿用\Term{monodromy-theorem}{单值性定理}算术证明的方法。第二个证明（第 IX 讲 3.6）则建立在对 \Term{neron-model}{Néron 模型}及其\Term{polarization}{极化}性质精细得多的分析之上。

Grothendieck 的第 I 至 V 讲未写成讲稿。这里将它们概括为第 I 讲。其中所陈述的结果都给出了简短而实质上完整的证明。

第 II 讲把 \Term{lefschetz-pencil}{Lefschetz 铅笔}（pinceaux de Lefschetz）方法和第 I 讲 5.3 用于\Term{fundamental-group}{基本群}的研究。对于\Term{characteristic-exponent}{特征指数}为 $p$ 的\Term{algebraically-closed-field}{代数闭域} $k$ 上的\Term{surface}{曲面} $S$，证明了\Term{fundamental-group}{基本群}的素于 $p$ 的\Term{completion}{完备化}——预\Term{finite-group}{有限群} $\pi_1^{(p)}(S,s)$——具有\Term{finite}{有限}的 pro-$(p)$ 展示。

如上所述，第 III 至 V 讲并不存在。

\SourcePage{8}
第 VI 讲叙述 Schlessinger 的\Term{deformation}{形变}理论，并作若干补充。其中注意计入所分类\Term{object}{对象}的\Term{infinitesimal-automorphism}{无穷小自同构}，避免过于草率地转到\Term{object}{对象}\Term{isomorphism}{同构}类的\Term{functor}{函子}。还给出了 $\operatorname{Ext}^i(L_{X/S},-)$ 的一种新构造，其中 $L_{X/S}$ 是 $X/S$ 的相对\Term{cotangent-complex}{余切复形}。在本讨论班其余各讲中，本讲的作用主要是其中对\Term{ordinary-quadratic-singularity}{普通二次奇点}\Term{deformation}{形变}的研究。

第 VII、VIII 讲专论\Term{biextension}{双扩张}理论。此理论在第 IX 讲中起着关键作用，用来表达从\Term{abelian-variety}{阿贝尔簇}转到其 \Term{neron-model}{Néron 模型}时\Term{polarization}{极化}如何变化。第 IX 讲包含\Term{abelian-variety}{阿贝尔簇}\Term{stable-reduction}{稳定约化}\Term{theorem}{定理}的证明及其若干应用。

本讨论班的续册 SGA 7 II，由 P. Deligne 和 N. Katz 编写，将于以后出版。

\begin{flushright}Bures-sur-Yvette，1972 年 5 月\\P. Deligne\end{flushright}

\SourcePage{9}
% 原书目录在本译本中由构建系统按中文标题及译本页码生成。
\SourcePage{10}
\SGAchapter{I}{A. Grothendieck 最初几讲的概要}{Résumé des premiers exposés de A. Grothendieck}
\begin{center}P. Deligne 撰写\end{center}
\textbf{本讲提要}
\begin{enumerate}
\item[0.] 预备知识
\item[1.] \Term{monodromy-theorem}{单值性定理}的算术证明
\item[2.] \Term{vanishing-cycles}{消失圈}
\item[3.] \Term{monodromy-theorem}{单值性定理}的几何证明
\item[4.] \Term{vanishing-cycles-sheaf}{消失圈层}的消没判据
\item[5.] \Term{monodromy}{单值性}对 $\pi_1$ 的\Term{action}{作用}
\item[6.] P. Deligne 的附录：\Term{stable-reduction}{稳定约化}\Term{theorem}{定理}的算术证明
\end{enumerate}

\SGAsection{0}{预备知识}
\textbf{0.0. 术语。}回顾下列定义。

\textbf{(0.0.1)} \Term{trait}{离散赋值环谱}（trait）：\Term{discrete-valuation-ring}{离散赋值环}的谱。通常用 $\eta$ 和 $s$ 分别表示它的\Term{generic-point}{一般点}和\Term{closed-point}{闭点}。

\textbf{(0.0.2)} \Term{strictly-local}{严格局部}：对环而言，指\Term{residue-field}{剩余域}\Term{separably-closed}{可分闭}的 \Term{henselian-local-ring}{Hensel 局部环}；对\Term{scheme}{概形}而言，指这种环的谱。

\textbf{(0.0.3)} $S$ 的\Term{geometric-point}{几何点}（在本讲中）：一个\Term{morphism}{态射} $\bar x:\operatorname{Spec}(\bar k)\to S$，其像为 $x\in S$，且 $k(\bar x)\overset{\mathrm{def}}=\bar k$ 是 $k(x)$ 的一个\Term{separable-closure}{可分闭包}。对于 \Term{henselian-trait}{Hensel 离散赋值环谱} $S=\operatorname{Spec}(V)$，常用 $\bar\eta$ 表示 $S$ 的一个\Term{geometric-generic-point}{几何一般点}（即像为 $\eta$ 的\Term{geometric-point}{几何点}）。\Term{valuation-ring}{赋值环} $\bar V$ 是 $V$ 在 $k(\bar\eta)$ 中的\Term{integral-closure}{整闭包}；它的\Term{residue-field}{剩余域}是 $k(s)$ 的一个\Term{separably-closed}{可分闭}\Term{algebra}{代数}\Term{group-extension}{扩张}，并确定 $S$ 的一个位于 $s$ 的\Term{geometric-point}{几何点} $\bar s$。

\SourcePage{11}
\textbf{(0.0.4)} $X$ 在\Term{geometric-point}{几何点} $\bar x$ 处的\Term{strict-localization}{严格局部化}：见 SGA 4 VIII 4.4。

\textbf{(0.0.5)} 在 $S$ 上\Term{essentially-finite-type}{本质有限型}：对\Term{local-scheme}{局部}（分别为 Hensel \Term{local-scheme}{局部}、\Term{strictly-local}{严格局部}）$S$ \Term{scheme}{概形}而言，指 $S$ 上某个\Term{finite-type}{有限型}\Term{scheme}{概形}的\Term{local-ring}{局部环}（分别为某\Term{local-ring}{局部环}的 \Term{henselization}{Hensel 化}、\Term{strict-localization}{严格局部化}）的谱。

\textbf{(0.0.6)} $F\mapsto F(n)$：\Term{tate-twist}{Tate 扭转}。

\textbf{(0.0.7)} \Term{finite-dimensional-vector-space}{有限维向量空间}的\Term{endomorphism}{自同态} $T$ 称为\Term{quasi-unipotent}{拟幺幂}的，若其\Term{eigenvalue}{特征值}都是\Term{root-of-unity}{单位根}；\Term{equivalence}{等价}地，存在 $N\ge1$、$M\ge1$，使 $(T^N-1)^M=0$。

\textbf{0.1. 关于 $\ell$ 进\Term{cohomology}{上同调}的回顾。}
设 $X$ 为\Term{characteristic-exponent}{特征指数}为 $p$ 的\Term{algebraically-closed-field}{代数闭域} $k$ 上的\Term{finite-type}{有限型}\Term{scheme}{概形}，$\ell$ 为与 $p$ 互素的\Term{prime-number}{素数}。

$X$ 的 $\ell$ 进\Term{cohomology}{上同调}群已在 SGA 4 和 SGA 5 中定义：
\begin{enumerate}
\item[a)] $H^i(X,\mathbf Z/\ell^n)$ 是 $X$ 的\Term{etale}{平展}\Term{site}{景} $X_{\text{ét}}$ 上以 $\mathbf Z/\ell^n$ 为系数的第 $i$ 个\Term{cohomology}{上同调}群；
\item[b)] 按定义，$H^i(X,\mathbf Z_\ell)=\varprojlim_n H^i(X,\mathbf Z/\ell^n)$；
\item[c)] 按定义，$H^i(X,\mathbf Q_\ell)=H^i(X,\mathbf Z_\ell)\otimes_{\mathbf Z_\ell}\mathbf Q_\ell$。
\end{enumerate}
$X$ 的\Term{proper-support}{固有支集} $\ell$ 进\Term{cohomology}{上同调}群同样由 SGA 4 XVII 中的 $H_c^i(X,\mathbf Z/\ell^n)$ 定义。

在下列各情形中，已能证明这些群 $H^\bullet(X)$ 具有适当的\Term{finiteness}{有限性}性质：

\textbf{(0.1.1)} 对\Term{proper-support}{固有支集}\Term{cohomology}{上同调}而言；

\textbf{(0.1.2)} $X$ 在 $k$ 上\Term{proper}{固有}；

\textbf{(0.1.3)} $X$ 在 $k$ 上\Term{smooth}{光滑}（由 a) 和 Poincaré \Term{duality}{对偶} SGA 4 XVIII）；

\textbf{(0.1.4)} 对 $k$ 上\Term{dimension}{维数} $\le\dim(X)$ 的\Term{finite-type}{有限型}\Term{scheme}{概形}有 Hironaka 式\Term{singularity-resolution}{奇点解消}；

\textbf{(0.1.5)} $i=0,1$。

\SourcePage{12}
在所有这些情形中，$H^i(X,\mathbf Z/\ell^n)$ 都是\Term{finite-group}{有限群}，$H^i(X,\mathbf Z_\ell)$ 都是\Term{finite-generated}{有限生成}模，$H^i(X,\mathbf Q_\ell)$ 都是\Term{finite-dimensional}{有限维}\Term{space}{空间}，而且（也许除 (0.1.5) 外）有\Term{split}{分裂}\Term{short-exact-sequence}{短正合列}
\[
0\longrightarrow H^i(X,\mathbf Z_\ell)\otimes\mathbf Z/\ell^n
\longrightarrow H^i(X,\mathbf Z/\ell^n)
\longrightarrow\operatorname{Tor}_1(H^{i+1}(X,\mathbf Z_\ell),\mathbf Z/\ell^n)
\longrightarrow0.
\]
此外，\Term{finiteness-theorem}{有限性定理}的证明每次都同时给出一个\Term{generic}{一般}\Term{constructibility-theorem}{可构造性定理}的证明：若 $k$ 是整的 Noether \Term{scheme}{概形} $S$ 的\Term{fraction-field}{分式域}，而 $X$ 是\Term{finite-type}{有限型}\Term{morphism}{态射} $f_1:X_1\to S$ 的\Term{generic-fibre}{一般纤维}，则存在\Term{nonempty}{非空}\Term{open-subset}{开集} $U\subset S$，使 $U$ 上的 $R^if_*\mathbf Z/\ell^n$ \Term{locally-constant}{局部常值}，且其形成与任意\Term{base-change}{基变换}相容。

\textbf{0.2.} 现在不再假设 $k$ 为\Term{algebraically-closed-field}{代数闭域}，令 $\bar k$ 为 $k$ 的一个\Term{algebraic-closure}{代数闭包}（或\Term{separable-closure}{可分闭包}，这里没有区别）。我们考虑“几何”\Term{cohomology}{上同调}群 $H^i(X_{\bar k})$，其中 $X_{\bar k}$ 是 $X$ 从 $k$ 到 $\bar k$ 作\Term{scalar-extension}{标量扩张}所得的\Term{scheme}{概形}。由结构搬移，\Term{galois-group}{Galois 群} $\operatorname{Gal}(\bar k/k)$ \Term{action}{作用}于这些群。在 $\mathbf Q_\ell$ \Term{cohomology}{上同调}中，若处于 (0.1.1) 至 (0.1.5) 的某种情形，就得到 $\operatorname{Gal}(\bar k/k)$ 在一个\Term{linear-group}{线性群} $\operatorname{GL}(n,\mathbf Q_\ell)$ 中的\Term{continuous-representation}{连续表示}。

\textbf{0.3. 关于\Term{trait}{离散赋值环谱}的回顾（0.0.1）。}
证明参见 [4]。设 $S$ 为 \Term{henselian-trait}{Hensel 离散赋值环谱}，$\eta,s,V,\bar V,\bar\eta,\bar s$ 如 (0.0.1)、(0.0.3)。以 $p$ 表示 $k(s)$ 的\Term{characteristic-exponent}{特征指数}，以 $\operatorname{Gal}(\bar\eta/\eta)$（分别以 $\operatorname{Gal}(\bar s/s)$）表示 \Term{galois-group}{Galois 群} $\operatorname{Gal}(k(\bar\eta)/k(\eta))$（分别为 $\operatorname{Gal}(k(\bar s)/k(s))$）。由结构搬移，$\operatorname{Gal}(\bar\eta/\eta)$ \Term{action}{作用}于 $k(\bar s)$，于是有从 $\operatorname{Gal}(\bar\eta/\eta)$ 到 $\operatorname{Gal}(\bar s/s)$ 的映射（实际上是\Term{surjective}{满射}），其核为\Term{inertia-group}{惯性群} $I$。$I$ 所确定的 $k(\bar\eta)$ 的子域是 $k(\eta)$ 的极大\Term{unramified}{非分歧}\Term{group-extension}{扩张}。

预\Term{finite-group}{有限群} $I$ 有一个最大的 pro-$p$ \Term{subgroup}{子群} $P$。\Term{quotient-group}{商群} $I/P$ 是\Term{tame-inertia-group}{温和惯性群}。\Term{canonical}{典范}地有
\[
I/P\simeq\varprojlim_n\mu_n(k(\bar s))\overset{\mathrm{def}}=\widehat{\mathbf Z}(1)(k(\bar s)).
\]
\SourcePage{13}
若 $\phi_n$ 为 $I/P$ 到 $\mu_n(\bar k)$ 的\Term{canonical-map}{典范映射}，$u$ 为 $\bar V$ 中\Term{valuation}{赋值}为 $1/n$ 的元素，则对 $\sigma\in I$ 有
\[
\sigma(u)\equiv\phi_n(u)u\pmod{\text{赋值大于 }1/n\text{ 的元素}}.
\]
综上，有
\[
\begin{gathered}
\operatorname{Gal}(\bar\eta/\eta)\supset I\supset P,\\
\operatorname{Gal}(\bar\eta/\eta)/I\simeq\operatorname{Gal}(\bar s/s),\\
I/P\simeq\widehat{\mathbf Z}(1)(k(\bar s)),\\
P:\ \text{pro-}p\text{ 群}\quad(p=1\text{ 时为 }\{e\}),
\end{gathered}\tag{0.3.1}
\]
且 $\operatorname{Gal}(\bar\eta/\eta)/I$ 在 $I/P$ 上由 $\operatorname{Gal}(\bar\eta/\eta)$ 中的\Term{inner-automorphism}{内自同构}给出的\Term{action}{作用}，可与 $\operatorname{Gal}(\bar s/s)$ 在 $\widehat{\mathbf Z}(1)(k(\bar s))$ 上的自然\Term{action}{作用}相等同。

\textbf{0.4.} 更\Term{generic}{一般}地，设 $S$ \Term{strictly-local}{严格局部}且\Term{regular}{正则}，$D$ 为 $S$ 中的\Term{regular-divisor}{正则除子}，$p'$ 为 $D$ 的\Term{generic-point}{一般点}的\Term{characteristic-exponent}{特征指数}。设 $\bar\eta$ 为 $S$ 的一个\Term{geometric-generic-point}{几何一般点}，$\bar s$ 为由它\Term{derive}{导出}的、位于\Term{closed-point}{闭点} $s$ 的\Term{geometric-point}{几何点}。由 Abhyankar \Term{lemma}{引理}，\Term{fundamental-group}{基本群} $\pi_1(S-D,\bar\eta)$ 有类似于 (0.3.1) 的\Term{resolution}{分解}：
\[
\begin{gathered}
\pi_1(S-D,\bar\eta)\supset I\supset P,\\
\pi_1(S-D,\bar\eta)/I\simeq\operatorname{Gal}(\bar s/s),\\
I/P\simeq\widehat{\mathbf Z}(1)(k(\bar s)),\\
P:\ \text{没有阶素于 }p'\text{ 的商群}.
\end{gathered}\tag{0.4.1}
\]
$\pi_1(S-D,\bar\eta)/I$ 在 $I/P$ 上由\Term{inner-automorphism}{内自同构}给出的\Term{action}{作用}，可与 $\operatorname{Gal}(\bar s/s)$ 在 $\widehat{\mathbf Z}(1)(k(\bar s))$ 上的自然\Term{action}{作用}相等同。

\textbf{0.5. Néron 解奇法。}
Néron 解奇法给出下列结果。
\begin{SGAresult}{\Term{lemma}{引理}}{0.5.1}
设 $S\to S_0$ 为 \Term{henselian-trait}{Hensel 离散赋值环谱}间的\Term{morphism}{态射}，其中 $S=\operatorname{Spec}(V)$、$S_0=\operatorname{Spec}(V_0)$，$\pi$ 为 $V_0$ 的一个\Term{uniformizer}{一致化元}。假设 $\pi$ 也是 $V$ 的\Term{uniformizer}{一致化元}，且\Term{group-extension}{扩张} $k(s)/k(s_0)$ 和 $k(\eta)/k(\eta_0)$ 都\Term{separable}{可分}。则 $V$ 是若干在 $V_0$ 上\Term{smooth}{光滑}、\Term{essentially-finite-type}{本质有限型}的 Hensel $V_0$ \Term{algebra}{代数} $B_i$ 的\Term{inductive-limit}{归纳极限}。
\end{SGAresult}
\SourcePage{14}
为构造所需的 $V_0$ \Term{algebra}{代数}，取 $V_0\subset B\subset V$，其中 $B$ 为\Term{finite-type}{有限型}，然后对 $\operatorname{Spec}(B)$ 施行 Néron 解奇法（[1] 第 4 节）。

设 $S=\operatorname{Spec}(V)$ 为 \Term{henselian-trait}{Hensel 离散赋值环谱}。

\textbf{(0.5.2)} 若 $S$ 为\Term{equal-characteristic}{等特征}，\Term{uniformizer}{一致化元}为 $\pi$，则可对 $S_0=\mathbf F_p[\pi]_{(\pi)}$ 应用 (0.5.1)。\Term{group-extension}{扩张} $s/s_0$ \Term{separable}{可分}，因为 $\mathbf F_p$ 为\Term{perfect-field}{完美域}；\Term{group-extension}{扩张} $\eta/\eta_0$ 也\Term{separable}{可分}，因为\Term{uniformizer}{一致化元} $\pi$ 不是 $p$ 次幂（当 $p\ne1$ 时），故可作为某个 $p$ 基的一部分。于是 $S$ 是若干\Term{smooth}{光滑}、\Term{essentially-finite-type}{本质有限型} Hensel $\mathbf F_p[\pi]_{(\pi)}$ \Term{algebra}{代数}的谱的\Term{limit}{极限}。

\textbf{(0.5.3)} 若 $S$ \Term{complete}{完备}、\Term{mixed-characteristic}{混合特征}，且其\Term{residue-field}{剩余域} $k$ \Term{perfect}{完美}，则 $V$ 为 \Term{witt-vectors}{Witt 向量}环（即 Cohen 环）$W(k)$ 的 Eisenstein \Term{group-extension}{扩张}：
\[
V\simeq W(k)[\pi]\big/\left(\pi^n+\sum_{i=0}^{n-1}a_i\pi^i\right).
\]
稍微扰动 $a_i$，会得到\Term{isomorphism}{同构}的\Term{group-extension}{扩张}。因此可假设 $a_i\in W(k')$，其中 $k'$ 在 $\mathbf F_p$ 上\Term{finite-generated}{有限生成}。令 $k_0$ 为 $k'$ 在 $k$ 中的\Term{perfect-closure}{完美闭包}，并令
\[
V_0=W(k_0)[\pi]\big/\left(\pi^n+\sum_{i=0}^{n-1}a_i\pi^i\right).
\]
\Term{lemma}{引理} 0.5.1 适用于 $V/V_0$，而 $V_0$ 的\Term{residue-field}{剩余域}是 $\mathbf F_p$ 上某个\Term{finite-generated}{有限生成}域的\Term{radicial}{根性}\Term{group-extension}{扩张}。

\SGAsection{1}{单值性定理的算术证明}
读者可在 [5] 的附录中找到 A. Grothendieck 的下列结果的证明。
\begin{SGAresult}{命题}{1.1}
设 $S$ 为 \Term{henselian-trait}{Hensel 离散赋值环谱}，$\ell$ 为与 $k(s)$ 的\Term{characteristic-exponent}{特征指数} $p$ 互素的\Term{prime-number}{素数}，$\rho$ 为 $\operatorname{Gal}(\bar\eta/\eta)$ 在 $\operatorname{GL}(n,\mathbf Q_\ell)$ 中的\Term{continuous-representation}{连续表示}。假设下列条件成立：
\[
(\star_\ell)\qquad
\text{$k(s)$ 的任何有限扩张都不包含所有阶为 $\ell$ 的幂的单位根。}
\]
\SourcePage{15}
则\Term{inertia-group}{惯性群} $I$ 有一个\Term{open-subgroup}{开子群} $I_1$，使得对 $\sigma\in I_1$，$\rho(\sigma)$ 都是\Term{unipotent}{幺幂}的。
\end{SGAresult}
若 $k(s)$ 在\Term{prime-field}{素域}上\Term{finite-generated}{有限生成}（或是这种域的\Term{radicial}{根性}\Term{group-extension}{扩张}），则条件 $(\star_\ell)$ 自动成立。

下列\Term{theorem}{定理}立即由 1.1 得到。
\begin{SGAresult}{\Term{monodromy-theorem}{单值性定理}}{1.2}
沿用上述记号，设 $X$ 为 $\eta$ 上的\Term{finite-type}{有限型}\Term{scheme}{概形}，$H$ 为 $X_{\bar\eta}$ 的以 $\mathbf Q_\ell$ 为系数、类型属于 (1.1.1) 至 (1.1.5) 之一的一个\Term{cohomology}{上同调}\Term{space}{空间}。若条件 $(\star_\ell)$ 成立，则 $I$ 的任意元素在 $H$ 上的\Term{action}{作用}都是\Term{quasi-unipotent}{拟幺幂}的。
\end{SGAresult}
0.5 中的\Term{limit}{极限}过渡结果可使我们去掉假设 $(\star_\ell)$。
\begin{SGAresult}{变式}{1.3}
1.2 中的条件 $(\star_\ell)$ 是多余的。
\end{SGAresult}
设 $H'_0$ 为产生 $H$ 的 $X_{\bar\eta}$ 的一个 $\mathbf Z_\ell$ \Term{cohomology}{上同调}\Term{space}{空间}，令 $H_0=H'_0/\text{挠子模}$。由假设，$H=H_0\otimes_{\mathbf Z_\ell}\mathbf Q_\ell$，且 $\operatorname{Gal}(\bar\eta/\eta)$ 的\Term{action}{作用}由下列表示\Term{derive}{导出}：
\[
\rho:\operatorname{Gal}(\bar\eta/\eta)\longrightarrow\operatorname{GL}(H_0).
\]
令 $I_{\ell'}$ 为 $I$ 的最大的素于 $\ell$ 的\Term{subgroup}{子群}。它在 $\operatorname{GL}(H_0)$ 中的像为素于 $\ell$ 的 $\ell$ 进 Lie 群，因而为\Term{finite-group}{有限群}。

作\Term{trait}{离散赋值环谱}的\Term{group-extension}{扩张} $S'\to S$。这不改变 $H_0$（SGA 4 XVI 6.6），且 $\rho(I')$ 在 $\rho(I)$ 中的指数\Term{finite}{有限}：这是因为 $I'/P'$ 在 $I/P$ 中的指数\Term{finite}{有限}，并结合上面的论述。因此，只需在 $S'$ 上证明\Term{theorem}{定理}；可以假设：
\begin{enumerate}
\item[a)] $\operatorname{Gal}(\bar\eta/\eta)$ 在 $H_0/\ell H_0$ 上的\Term{action}{作用}平凡；
\item[b)] 若 $S$ 为\Term{mixed-characteristic}{混合特征}，则 $S$ \Term{complete}{完备}且\Term{residue-field}{剩余域}\Term{perfect}{完美}（EGA $0_{\mathrm{III}}$ 10.3.1）。
\end{enumerate}
令 $S=\operatorname{Spec}(V)$，如 (0.5.2)、(0.5.3) 那样定义 $V_0$，并应用 (0.5.1)。有
\[
V=\varinjlim_i B_i,
\]
\SourcePage{16}
其中 $S_i=\operatorname{Spec}(B_i)$ 在 $V_0$ 上本质\Term{smooth}{光滑}且为\Term{finite-type}{有限型}。设 $D_i\subset S_i$ 由方程 $\pi=0$ 定义。当 $i$ 充分大时，由“\Term{generic}{一般}\Term{constructibility}{可构造性}”（0.1），$H_0$ 来自 $S_i-D_i$ 上的一个\Term{twisted-constant}{扭常值} $\mathbf Z_\ell$ \Term{sheaf}{层} $\mathscr H_i$，且 $\mathscr H_i/\ell\mathscr H_i$ 为\Term{constant-sheaf}{常值层}。

\Term{resolution}{分解} (0.3.1) 和 (0.4.1) 给出\Term{commutative-diagram}{交换图}：
\[
\begin{tikzcd}[column sep=small]
\operatorname{Gal}(\bar\eta/\eta)\arrow[d]\arrow[r,phantom,"\supset"]&I\arrow[d]\arrow[r,phantom,"\supset"]&P\arrow[d]\\
\pi_1(S_i-D_i,\bar\eta)\arrow[r,phantom,"\supset"]&I_i\arrow[r,phantom,"\supset"]&P_i
\end{tikzcd}
\qquad
\begin{tikzcd}
I/P\arrow[d]\arrow[r,"\sim"]&\widehat{\mathbf Z}(1)\arrow[d,equal]\\
I_i/P_i\arrow[r,"\sim"]&\widehat{\mathbf Z}(1)
\end{tikzcd}\tag{1.3.1}
\]
$\operatorname{Gal}(\bar\eta/\eta)$ 的表示 $H_0$ 由 $\pi_1(S_i-D_i,\bar\eta)$ 在 $H_0$ 上的一个\Term{action}{作用}\Term{derive}{导出}，后者在 $H_0/\ell H_0$ 上平凡。群 $\operatorname{Ker}(\operatorname{GL}(H)\to\operatorname{GL}(H/\ell H))$ 是 pro-$\ell$ 群；故 $P_i$ 的\Term{action}{作用}平凡。再应用下述 1.1 的变式。

\begin{SGAresult}{变式}{1.4}
沿用 0.4 的记号，设 $\rho$ 为 $\operatorname{Gal}(\bar\eta/\eta)$ 在 $\operatorname{GL}(n,\mathbf Q_\ell)$ 中的\Term{continuous-representation}{连续表示}。假设 $k(s)$ 满足 $(\star_\ell)$，且 $\rho(P)$ \Term{finite}{有限}。则 $I$ 的\Term{action}{作用}是\Term{quasi-unipotent}{拟幺幂}的。
\end{SGAresult}

\SGAsection{2}{消失圈}
\textbf{2.1.} 设 $S$ 为\Term{strictly-local}{严格局部}\Term{trait}{离散赋值环谱}，沿用 0.3 的记号。由假设，这里 $s=\bar s$。设 $f:X\to S$ 为 $S$ 上的\Term{finite-type}{有限型}\Term{scheme}{概形}。\Term{vanishing-cycles}{消失圈}\Term{form}{形式}体系是一种工具：从关于\Term{morphism}{态射} $f$ 的\Term{local-scheme}{局部}信息（在 $X$ 上）出发，研究 $X_s$ 与 $X_{\bar\eta}$ 的\Term{cohomology}{上同调}之间的关系，以及\Term{inertia-group}{惯性群}在后者上的\Term{action}{作用}。

这类结果中的一个关键是\Term{specialization}{特化}\Term{theorem}{定理} SGA 4 XVI 2.2：若 $f$ \Term{proper}{固有}且\Term{smooth}{光滑}，则 $X_s$ 与 $X_{\bar\eta}$ 有相同的\Term{cohomology}{上同调}。

在\Term{transcendental}{超越}理论中，若 $f:X\to D$ 为\Term{analytic-space}{解析空间} $X$ 到\Term{disc}{圆盘} $D$ 的\Term{proper}{固有}\Term{morphism}{态射}，\Term{vanishing-cycles}{消失圈}如下出现：

\textbf{a)} 缩小 $D$ 后，可假设 $X$ 在去心\Term{disc}{圆盘} $D^*$ 上是\Term{topology}{拓扑}\Term{fibre-bundle}{纤维丛}（Thom \Term{isotopy}{同痕}\Term{theorem}{定理}）。特别地，所有\Term{fibre}{纤维} $X_t$（$t\in D^*$）都\Term{isomorphism}{同构}。
\SourcePage{17}
\textbf{b)} 把\Term{special-fibre}{特殊纤维} $X_s$ 想象为由“那个”\Term{generic-fibre}{一般纤维} $X_t$ \Term{retraction}{收缩}其中某些多面体而得到，是合理的；这样就会有 $\pi:X_t\to X_s$。于是，\Term{vanishing-cycles}{消失圈}方法就是借助 $\pi$ 的 \Term{leray-spectral-sequence}{Leray 谱序列}来研究 $X_s$ 与 $X_t$ 的\Term{cohomology}{上同调}之差。

下面说明的几何理论与上述\Term{transcendental}{超越}理论很接近；应以 $S$ 的\Term{geometric-generic-point}{几何一般点} $\bar\eta$ 代替\Term{generic-point}{一般点} $t\in D$。反过来，也可按照几何理论来建立\Term{transcendental}{超越}理论：以 $D^*$ 的\Term{universal-cover}{万有覆盖} $\widetilde D^*$ 代替 $t$ 或 $\bar\eta$，以 $X\times_D\widetilde D^*$ 代替 $X_t$ 或 $X_{\bar\eta}$。这样可去掉许多 $\varepsilon$ 和 $\eta$，但有时也会丢失信息。

\textbf{2.2.} 固定一个\Term{torsion-ring}{挠环} $\Lambda$，使 $p$ 在其中\Term{invertible}{可逆}（例如 $\Lambda=\mathbf Z/\ell^n\mathbf Z$，其中 $\ell$ 素于 $p$）。考虑\Term{morphism}{态射}图
\[
\begin{tikzcd}
X_s\arrow[r,"i"]\arrow[d,"f_s"']&X\arrow[d,"f"]&&X_{\bar\eta}\arrow[ll,"\bar j"']\arrow[d,"f_{\bar\eta}"]\\
s\arrow[r]&S&\eta\arrow[l]&\bar\eta\arrow[l]
\end{tikzcd}
\]
$X_s$ 上的\Term{vanishing-cycles}{消失圈}\Term{sheaf}{层} $\psi^i$ 为
\[
\psi^i=i^*R^i\bar j_*\Lambda.\tag{2.2.1}
\]
若 $f$ \Term{proper}{固有}，则有（SGA 4 XII 5.5）
\[
i^*:H^p(X,R^q\bar j_*\Lambda)\xrightarrow{\sim}
H^p(X_s,i^*R^q\bar j_*\Lambda),\tag{2.2.2}
\]
且 $\bar j$ 的 \Term{leray-spectral-sequence}{Leray 谱序列}可改写为
\[
H^p(X_s,\psi^q)\Longrightarrow H^{p+q}(X_{\bar\eta},\Lambda).\tag{2.2.3}
\]
当无穷远处没有变化时，有时不加\Term{properness}{固有性}假设也可得到这个\Term{spectral-sequence}{谱序列}。无穷远处确实需要某种假设，这从平凡情形 $X_s=\varnothing$ 就能看出。

\SourcePage{18}
由定义立即得到：
\begin{SGAresult}{命题}{2.3}
设 $\bar x$ 为 $X_s$ 的一个\Term{geometric-point}{几何点}（例如一个\Term{closed-point}{闭点}），$X_{(\bar x)}$ 为 $X$ 在 $\bar x$ 处的严格 \Term{henselization}{Hensel 化}，$X_{(\bar x)\bar\eta}$ 为 $X_{(\bar x)}$ 到 $S$ 的\Term{projective-object}{投射}的\Term{geometric-generic-fibre}{几何一般纤维}。则
\[
\psi^i_{\bar x}=H^i(X_{(\bar x)\bar\eta},\Lambda).
\]
\end{SGAresult}
\begin{SGAresult}{推论}{2.4}
在 $f$ \Term{smooth}{光滑}之处，$i>0$ 时 $\psi^i=0$，且 $\psi^0=\Lambda$。
\end{SGAresult}
这就是 SGA 4 XV 2.1。

群 $\operatorname{Gal}(\bar\eta/\eta)$（由假设，它等于\Term{inertia-group}{惯性群}）通过结构搬移\Term{action}{作用}于 $\psi^i$，并且：
\begin{SGAresult}{附论}{2.5}
\Term{spectral-sequence}{谱序列} 2.2.3 是 $\operatorname{Gal}(\bar\eta/\eta)$ \Term{equivariant}{等变}的。
\end{SGAresult}
\begin{SGAresult}{变式}{2.6}
设 $\bar S=\operatorname{Spec}(\bar V)$、$\bar X=X\times_S\bar S$。$X_{\bar\eta}$ 是 $X_s=X_{\bar s}$ 在 $\bar X$ 中的补集。令
\[
\phi^i=\mathcal H_{X_s}^{i+1}(\bar X,\Lambda)
\]
（\Term{support}{支集}\Term{cohomology}{上同调}\Term{sheaf}{层}）。有\Term{sheaf}{层}的\Term{exact-sequence}{正合列}
\[
0\longrightarrow\phi^{-1}\longrightarrow\Lambda\longrightarrow\psi^0
\longrightarrow\phi^0\longrightarrow0,
\]
且当 $i>0$ 时有\Term{isomorphism}{同构} $\psi^i\xrightarrow{\sim}\phi^i$。若 $f$ \Term{smooth}{光滑}，则所有 $\phi^i$ 都为零。

令 $\phi^i_{\mathrm{gl}}=H^i_{X_s}(\bar X,\Lambda)$（$\phi^i$ 的整体对应物）。有\Term{spectral-sequence}{谱序列}
\[
H^p(X_s,\phi^q)\Longrightarrow\phi^{p+q}_{\mathrm{gl}},\tag{2.6.1}
\]
以及一个\Term{cohomology}{上同调}\Term{long-exact-sequence}{长正合列}；当 $f$ \Term{proper}{固有}时，它可改写为
\[
\cdots\longrightarrow H^i(X_s,\Lambda)\longrightarrow H^i(X_{\bar\eta},\Lambda)
\longrightarrow\phi^i_{\mathrm{gl}}\longrightarrow\cdots .\tag{2.6.2}
\]
\end{SGAresult}
\begin{SGAresult}{变式}{2.7}
设 $K_t$ 为 $K$ 在 $\bar K$ 中的极大\Term{tame-ramification}{温和分歧}\Term{group-extension}{扩张}。令 $X_t=X\times_S\operatorname{Spec}(K_t)$，并以 $\bar j_t$ 表示 $X_t$ 到 $X$ 的\Term{projective-object}{投射}。定义 $\psi^i$ 的温和变式（“tame”）如下：
\SourcePage{19}
\[
\psi_t^i=R^i\bar j_{t*}\Lambda.
\]
利用 $P$ 为 pro-$p$ 群且 $p$ 在 $\Lambda$ 中\Term{invertible}{可逆}，得到
\[
H^*(X_t,\Lambda)=H^*(X_{\bar\eta},\Lambda)^P,\tag{2.7.1}
\]
\[
\psi_t^i=(\psi^i)^P.\tag{2.7.2}
\]
若 $f$ \Term{proper}{固有}，则 $\bar j_t$ 的 \Term{leray-spectral-sequence}{Leray 谱序列}由 2.2.3 取 $P$ \Term{invariant}{不变}元得到：
\[
H^p(X_s,\psi_t^q)\Longrightarrow H^{p+q}(X_t,\Lambda)
=H^{p+q}(X_{\bar\eta},\Lambda)^P.\tag{2.7.3}
\]
\end{SGAresult}

\SGAsection{3}{单值性定理的几何证明}
设 $\Lambda$ 如 2.2。
\begin{SGAresult}{猜想（\Term{purity}{纯性}）}{3.1}
设 $A$ 为\Term{excellent}{优良}、\Term{strictly-local}{严格局部}（0.2）的\Term{regular}{正则}环，$D$ 为 $\operatorname{Spec}(A)$ 中由一个参数 $x$ 定义的\Term{regular-divisor}{正则除子}。令 $U=\operatorname{Spec}(A)-D=\operatorname{Spec}(A[1/x])$。则
\[
H^q(U,\Lambda)=
\begin{cases}
\Lambda,&q=0,\\
\Lambda(-1),&q=1\quad\text{（记号见 (0.0.6)）},\\
0,&q>1.
\end{cases}\tag{3.1.1}
\]
\end{SGAresult}
此猜想在以下情形已知成立（SGA 4 XIX）：$q=0$ 或 $1$；$A$ 为\Term{characteristic}{特征} $0$；$A$ 为\Term{equal-characteristic}{等特征} $p$，且在\Term{characteristic}{特征} $p$、\Term{dimension}{维数} $\le\dim(A)$ 时有\Term{singularity-resolution}{奇点解消}；或 $(\operatorname{Spec}(A),D)$ 为\Term{smooth}{光滑}对（相对\Term{purity}{纯性}\Term{theorem}{定理} SGA 4 XVI 3.7）。

当\Term{purity}{纯性}可用，且 $D=\sum_{i\in B}D_i$ 为 $\operatorname{Spec}(A)$ 中由\Term{regular}{正则}\Term{system-of-parameters}{参数系}的一部分定义的\Term{normal-crossings-divisor}{正规交叉除子}时，$U=\operatorname{Spec}(A)-D$ 的\Term{cohomology}{上同调}由下式给出：
\SourcePage{20}
\[
\begin{aligned}
H^0(U,\Lambda)&=\Lambda,\\
H^1(U,\Lambda)&=\Lambda(-1)^B,\\
H^q(U,\Lambda)&=\bigwedge^qH^1(U,\Lambda)\quad\text{（由杯积给出）}.
\end{aligned}\tag{3.1.2}
\]
这些公式与在 $\mathbf C$ 上给出 $(D^*)^B\times D^k$ 的\Term{cohomology}{上同调}的公式相同。

\textbf{3.2.} 设 $S$ 为\Term{strictly-local}{严格局部}\Term{trait}{离散赋值环谱}。保留 (0.0.1)、(0.0.3) 的记号。设 $f:X\to S$ 为\Term{finite-type}{有限型}\Term{morphism}{态射}。假设 $X$ \Term{regular}{正则}、$X_\eta$ \Term{smooth}{光滑}，且 $X_s$ 为 $X$ 中的\Term{normal-crossings-divisor}{正规交叉除子}。假设此\Term{divisor}{除子}是一些\Term{regular-divisor}{正则除子}的和（在整体上，而非仅对\Term{etale}{平展}\Term{topology}{拓扑}\Term{local-scheme}{局部}如此）：
\[
(X_s)_{\mathrm{red}}=\sum_{i\in B}D_i
\]
（\Term{generic}{一般}情形可由\Term{localization}{局部化}得到）。记这些\Term{divisor}{除子}的\Term{multiplicity}{重数}为 $m_i$：
\[
X_s=\sum_{i\in B}m_iD_i.
\]
设 $x$ 为 $X_s$ 的一点，$\bar x$ 为 $X$ 的一个位于 $x$ 的\Term{geometric-point}{几何点}（例如 $x$ 为\Term{closed-point}{闭点}，$\bar x=x$），$C=\{i\mid x\in D_i\}$。令 $K$ 为集中于\Term{degree}{次数} $0$ 和 $-1$ 的\Term{complex}{复形}
\[
K:\mathbf Z^C\xrightarrow{d}\mathbf Z,\qquad d((n_j))=\sum_jm_jn_j.
\]
\begin{SGAresult}{\Term{theorem}{定理}（使用\Term{purity}{纯性}）}{3.3}
温和\Term{vanishing-cycles}{消失圈}群 $\psi^i_{t\bar x}$ 如下：

\textup{(i)} 有\Term{graded}{分次} $\Lambda$ \Term{algebra}{代数}的\Term{canonical}{典范}\Term{isomorphism}{同构}
\[
\psi^*_{t\bar x}\simeq
\bigwedge\nolimits^*\bigl(H_1(K)\otimes\Lambda(-1)\bigr)
\otimes_\Lambda\psi^0_{t\bar x}.
\]
\textup{(ii)} $\psi^0_{t\bar x}$ 是 $\Lambda$ \Term{algebra}{代数} $\Lambda^E$，其中 $E$ 为一个\Term{set}{集合}，\Term{tame-inertia-group}{温和惯性群} $I_t=I/P=\widehat{\mathbf Z}(1)(k(s))$ 在 $E$ 上\Term{transitive-action}{传递作用}，且 $E$ 的元素个数为 $\#H_0(K)$ 的最大的素于 $p$ 的因子。
\end{SGAresult}
证明方法如下。

\textbf{a)} 设 $X_{(\bar x)}$ 为 $X$ 在 $\bar x$ 处的\Term{strict-localization}{严格局部化}，$U=X_{(\bar x)}-X_{(\bar x)s}$（一个\Term{normal-crossings-divisor}{正规交叉除子}的补集），$t_j$ 为 $D_j$ 的\Term{local-scheme}{局部}方程，令
\[
X_n=X_{(\bar x)}[(t_j^{1/n})]\qquad((n,p)=1;
\text{ 这是一个正则概形}),
\]
并令 $U_n$ 为 $U$ 在 $X_n$ 中的\Term{inverse-image}{逆像}。由 3.1.2，有
\[
H^q(U_n,\Lambda)\simeq\bigwedge^q\bigl(\Lambda(-1)^C\bigr).
\]
若 $m=nd$，则从 $H^q(U_n,\Lambda)$ 到 $H^q(U_m,\Lambda)$ 的“\Term{inverse-image}{逆像}”映射为乘以 $d^q$。令
\SourcePage{21}
\[
\widetilde U=\varprojlim_{(n,p)=1}U_n,
\]
于是
\[
\begin{cases}
H^0(\widetilde U,\Lambda)=\Lambda,\\
H^q(\widetilde U,\Lambda)=\varinjlim_nH^q(U_n,\Lambda)=0,&q\ne0.
\end{cases}\tag{3.3.1}
\]
\textbf{b)} $\widetilde U$ 是 $U$ 的一个\Term{covering}{覆盖}，其群为 $\widehat{\mathbf Z}(1)(k(s))^C$。由定义，有 $\psi^*_{t\bar x}=H^*(U_{\eta_t},\Lambda)$；$\widetilde U$ 又是 $U_{\eta_t}$ 的每个\Term{connected-component}{连通分支}的\Term{covering}{覆盖}，其群为 $H_1(K)\otimes\widehat{\mathbf Z}(1)(k(s))$，而\Term{tame-inertia-group}{温和惯性群}传递置换这些\Term{connected-component}{连通分支}。由 (3.3.1) 和 Hochschild–Serre \Term{spectral-sequence}{谱序列}，便得到 3.3 的公式。

形象地说，可以把 $U_{\bar\eta}$ 看成一个\Term{morphism}{态射} $\delta$ 的\Term{fibre}{纤维}；$\delta$ 从\Term{cohomology}{上同调}意义下的 $K(\widehat{\mathbf Z}(1)(k(s))^C,1)$ \Term{space}{空间} $U$ 映到\Term{cohomology}{上同调}意义下的 $K(\widehat{\mathbf Z}(1)(k(s)),1)$ \Term{space}{空间} $\eta$，并诱导
\[
d=\sum_jm_jn_j:\widehat{\mathbf Z}(1)(k(s))^C\longrightarrow
\widehat{\mathbf Z}(1)(k(s)).
\]
\begin{SGAresult}{推论（使用\Term{purity}{纯性}）}{3.4}
假设 $f$ \Term{proper}{固有}。设 $N$ 为各\Term{multiplicity}{重数} $m_j$（$j\in B$）的最小公倍数。设 $q\ge0$，$q'$ 为 $q$ 与经过同一点的\Term{divisor}{除子} $D_j$ 的最大个数两者的\Term{infimum}{下确界}。对\Term{tame-inertia-group}{温和惯性群} $I_t=I/P$ 中的 $T$，有
\[
(T^N-1)^{q'}=0\qquad\text{在 }H^q(X_{\bar\eta},\Lambda)^P\text{ 上}.
\]
\end{SGAresult}
这立即由 3.3 及 2.5 的温和变式 2.7 得出。

对 $H^1$、在\Term{characteristic}{特征} $0$ 时，或当 $S$ 在某个域上\Term{essentially-finite-type}{本质有限型}时，所需的\Term{purity}{纯性}结果都是可用的。

\begin{SGAresult}{\Term{theorem}{定理}}{3.5}
设 $C$ 为\Term{discrete-valuation-ring}{离散赋值环} $V$ 的\Term{fraction-field}{分式域} $k(\eta)$ 上的\Term{smooth}{光滑}\Term{projective}{射影}\Term{curve}{曲线}。\Term{inertia-group}{惯性群} $I$ 在 $H^1(C_{\bar\eta},\mathbf Q_\ell)$ 上通过二阶\Term{quasi-unipotent}{拟幺幂}\Term{automorphism}{自同构}\Term{action}{作用}：存在 $N$，使得对 $T\in I$，有
\[
(T^N-1)^2=0\qquad\text{在 }H^1(C_{\bar\eta},\mathbf Q_\ell)\text{ 上}.
\]
\end{SGAresult}
归结到 $V$ \Term{strictly-local}{严格局部}的情形。所要断言的是：$I$ 有一个\Term{open-subgroup}{开子群} $I_1$，使 $T\in I_1$ 时，$H^1$ 上有 $(T-1)^2=0$。
\SourcePage{22}
我们可以把标量从 $k(\eta)$ \Term{group-extension}{扩张}到它的一个\Term{finite-extension}{有限扩张}（以 $V$ 在此\Term{group-extension}{扩张}中的\Term{normalization}{正规化}代替 $V$）。由于 pro-$p$ 群 $P$ 在 $\ell$ 进 Lie 群 $\operatorname{GL}(n,\mathbf Q_\ell)$ 中的像为\Term{finite-group}{有限群}，特别可假设 $P$ 的\Term{action}{作用}平凡。

设 $C_0$ 为 $C$ 在 $V$ 上的\Term{minimal-model}{极小模型}（$C_0$ 的存在性基于 Abhyankar \Term{theorem}{定理}，见 [3] 第 2 节第 87 页的讨论；这里也可通过取 $V$ 的\Term{completion}{完备化}而归结到 $V$ \Term{excellent}{优良}的情形，使用 SGA 4 XVI 1.6）。设 $C_1$ 为由 $C_0$ 经\Term{blowup}{爆破}得到的 $C$ 的\Term{regular}{正则}\Term{model}{模型}，其\Term{special-fibre}{特殊纤维}为\Term{normal-crossings-divisor}{正规交叉除子}。3.3、3.4 的论证适用于 $C_1$、$\psi_t^0$、$\psi_t^1$ 及
\[
H^1(C_{\bar\eta},\mathbf Z/\ell^n)^P=H^1(C_{\bar\eta},\mathbf Z/\ell^n),
\]
其中 $N$ 与 $n$ 无关。3.5 因此得证。

\textbf{注记 3.6。} 设 $A$ 为 $k(\eta)$ 上的\Term{abelian-variety}{阿贝尔簇}。由于 $A$ 是某个 Jacobian 簇的商，仍可得到：$I$ 在 $H^1(A_{\bar\eta},\mathbf Q_\ell)$ 或 $T_\ell(A)$ 上的\Term{action}{作用}为二阶\Term{quasi-unipotent}{拟幺幂}。

\textbf{注记 3.7。} 事实上，对 $\eta$ 上任意\Term{finite-type}{有限型}\Term{scheme}{概形} $X$，$I$ 在 $H^1(X_{\bar\eta},\mathbf Q)$ 上的\Term{action}{作用}都是二阶\Term{quasi-unipotent}{拟幺幂}。

\Term{theorem}{定理} 3.5（或 3.6）是第 IX 讲第 3 节将给出的\Term{abelian-variety}{阿贝尔簇}\Term{stable-reduction}{稳定约化}\Term{theorem}{定理}证明的基础。

\textbf{3.8.} 显然，当\Term{purity}{纯性}和\Term{singularity-resolution}{奇点解消}可用时（例如在\Term{characteristic}{特征} $0$ 时），前面的方法证明了\Term{monodromy-theorem}{单值性定理}，并给出了\Term{nilpotence-index}{幂零指数}的界。例如，在\Term{equal-characteristic}{等特征} $0$ 时得到：若 $X$ 在 $\eta$ 上\Term{proper}{固有}且\Term{smooth}{光滑}，$T$ 属于 $I$ 的某个\Term{open-subgroup}{开子群}，则
\[
(T-1)^{q+1}=0\qquad\text{在 }H^q(X_{\bar\eta},\mathbf Z_\ell)\text{ 上}
\qquad(T\in I_1\subset I).
\]

\SGAsection{4}{消失圈层的消没判据}
\textbf{4.1.} 设 $S$ 为\Term{strictly-local}{严格局部}\Term{trait}{离散赋值环谱}，$f:X\to S$ 为 $S$ 上的\Term{finite-type}{有限型}\Term{scheme}{概形}。令 $n=\dim(X_\eta)$。

若 $X'$ 为 $X_\eta$ 在 $X$ 中的\Term{scheme}{概形}\Term{closure}{闭包}，则：

\textbf{a)} 在 $X_s-X'_s$ 上，$\phi^{-1}=\Lambda$，其余 $\phi^i$ 都为零；
\SourcePage{23}
\textbf{b)} 在 $X'_s$ 上，属于 $X/S$ 和 $X'/S$ 的 $\phi^i$、$\psi^i$ 分别相同，且 $\phi^{-1}=0$。

由于 $X'$ 在 $S$ 上\Term{flat}{平坦}，这往往可使问题归结到 $X$ 在 $S$ 上\Term{flat}{平坦}的情形。

\begin{SGAresult}{\Term{theorem}{定理}}{4.2}
\textup{(i)} 当 $i>n$ 时，$\psi^i=0$。

\textup{(ii)} 更确切地，设 $x$ 为 $X_s$ 的一点，其\Term{closure}{闭包}的\Term{dimension}{维数}为 $k$（$k=\operatorname{trdeg}(k(x)/k(s))$），$\bar x$ 为位于 $x$ 的\Term{geometric-point}{几何点}，则
\[
\psi^i_{\bar x}=0\quad(i>n-k),\qquad\text{即}\quad d(\psi^i)\le n-i
\]
（采用 SGA 4 XIV 2.1 的记号）。
\end{SGAresult}
有 $\psi^i_{\bar x}=H^i(X_{(\bar x)\bar\eta},\Lambda)$，而 $X_{(\bar x)\bar\eta}$ 是 $k(\bar\eta)$ 上\Term{dimension}{维数} $\le n$ 的\Term{affine}{仿射}\Term{scheme}{概形}的\Term{inverse-limit}{投射极限}。因此由\Term{affine}{仿射} Lefschetz \Term{theorem}{定理}（SGA 4 XIV 3.2）作\Term{limit}{极限}过渡即得 (i)。

为证明 (ii)，由 4.1 可假设 $X$ 在 $S$ 上\Term{flat}{平坦}，且存在\Term{quasi-finite}{拟有限}、\Term{flat}{平坦}的 $S$ \Term{morphism}{态射} $g:X\to\mathbf A_S^k$，使 $g(x)$ 为 $\mathbf A_s^k$ 的\Term{generic-point}{一般点}。设 $S'$ 为 $\mathbf A_S^k$ 在 $\bar x$ 处的\Term{strict-localization}{严格局部化}\Term{trait}{离散赋值环谱}。则：

\textbf{a)} $S'_{\bar\eta}$ 为一个域的谱，且 $\operatorname{Gal}(\bar\eta'/S'_{\bar\eta})$ 为某个 $p$ 群 $Q$。

\textbf{b)} 将 (i) 用于 $g$，有 $H^i(X_{(\bar x)\bar\eta'},\Lambda)=0$（$i>n-k$）。

\textbf{c)} 有
\[
H^i(X_{(\bar x)\bar\eta},\Lambda)
=H^i(X_{(\bar x)\bar\eta'},\Lambda)^Q,
\]
\Term{theorem}{定理}由此得证。

\begin{SGAresult}{推论}{4.3}
若 $f$ \Term{proper}{固有}且\Term{flat}{平坦}，并且 $f$ 在 $X_s$ 中的非\Term{smooth}{光滑}点集的\Term{dimension}{维数} $\le d$，则\Term{specialization}{特化}\Term{morphism}{态射}
\[
H^i(X_s)\longrightarrow H^i(X_{\bar\eta})
\]
在 $i>n+d+1$ 时为\Term{isomorphism}{同构}，在 $i=n+d+1$ 时为\Term{surjective}{满射}。
\end{SGAresult}
\textbf{4.4.} 为在某些情形证明低次 $\phi^i$ 为零，我们将借助\Term{local-scheme}{局部}\Term{duality}{对偶}理论（直接使用，或通过 SGA 2 XIV 的\Term{local-scheme}{局部} Lefschetz \Term{theorem}{定理}）。这一工具仅在\Term{characteristic}{特征} $0$ 时可用（或在\Term{equal-characteristic}{等特征} $p$ 时，对所涉及的\Term{dimension}{维数}有 Hironaka 式\Term{singularity-resolution}{奇点解消}时可用）。

\SourcePage{24}
\begin{SGAresult}{\Term{theorem}{定理}}{4.5}
假设 $S$ 为\Term{characteristic}{特征} $0$。设 $f:X\to S$ 为\Term{finite-type}{有限型}\Term{morphism}{态射}，$x$ 为 $X_s$ 的一个\Term{closed-point}{闭点}。假设：

\textup{(a)} $x$ 是其所在\Term{fibre}{纤维}中的孤立非\Term{smooth}{光滑}点；

\textup{(b)} $X$ 在 $x$ 处的相对\Term{geometric-depth}{几何深度} $\ge n$：在 $x$ 的某个\Term{neighborhood}{邻域}中，$X$ 可看作一个在 $x$ 处\Term{relative-dimension}{相对维数}为 $N$ 的\Term{smooth}{光滑}\Term{scheme}{概形}内由 $k$ 个方程定义的\Term{subscheme}{子概形}，并且 $n\le N-k$。

则 $i<n$ 时，$\phi_x^i=0$。
\end{SGAresult}
设 $X_x$ 为 $X$ 在 $x$ 处的\Term{strict-localization}{严格局部化}，$V=X_x-\{x\}$。对每个在 $S$ 上\Term{finite}{有限}的\Term{trait}{离散赋值环谱} $S'$，同样令 $X'_x=X_x\times_SS'$、$V'=V\times_SS'$。设 $\bar S$ 为 $S$ 在 $\bar\eta$ 中的\Term{normalization}{正规化}，令 $\bar X_x=X_x\times_S\bar S$、$\bar V=V\times_S\bar S$。有：

\textbf{(1)} 若当 $i\le m$ 时，$\phi^i=0$ 在 $(X_x)_s$ 上成立（即在 $x$ 于 $X_s$ 中的每个\Term{generalization}{一般化点}处成立），则
\[
H_{\{x\}}^{i-1}(\bar X_x)=\widetilde H^i(\bar V)
\xrightarrow{\sim}\widetilde H^i(\bar V_{\bar\eta})=\phi^i
\qquad(i\le m).
\]
事实上，$\phi^i$ 是 $\bar V$ 中以 $V_s$ 为\Term{support}{支集}的\Term{cohomology}{上同调}\Term{sheaf}{层}；在 4.5 的假设下，(a) 可对 $m=\infty$ 应用（SGA 4 XV 2.1）。

\textbf{(2)} 4.5 的假设 (b) 在\Term{base-change}{基变换}下\Term{stable}{稳定}。由 SGA 2 XIV 5.6，它蕴涵
\[
\widetilde H^i(V')=H^{i-1}_{\{x\}}(X'_x)=0\qquad(i<n).
\]
注意 $\widetilde H^i(\bar V)$ 是 $\widetilde H^i(V')$ 的\Term{inductive-limit}{归纳极限}，即得结论。

由 4.2 和 4.5，得到：
\begin{SGAresult}{推论}{4.6}
设 $S$ 为\Term{characteristic}{特征} $0$ 的 \Term{henselian-trait}{Hensel 离散赋值环谱}，$X/S$ 为\Term{relative-dimension}{相对维数} $n$ 的相对\Term{complete-intersection}{完全交}，$x\in X_s$ 为其所在\Term{fibre}{纤维}中的孤立非\Term{smooth}{光滑}点。则 $i\ne n$ 时，$\phi_x^i=0$。
\end{SGAresult}
\begin{SGAresult}{推论}{4.7}
在 4.6 的假设下，$\phi_x^n$ 为\Term{flat}{平坦} $\Lambda$ 模。
\end{SGAresult}
事实上，对\Term{finite}{有限} $\Lambda$ \Term{algebra}{代数} $\Lambda'$，\Term{generic}{一般}\Term{form}{形式}论证（SGA 4 XVII 5.2.11）对属于 $\Lambda'$ 的 $\phi_{\Lambda'}^i$ 给出
\SourcePage{25}
\[
\phi_{\Lambda'}^{n-i}=\operatorname{Tor}_i^\Lambda(\Lambda',\phi_\Lambda^n),
\]
然后对 $\Lambda'$ 应用 4.6 即可。

\begin{SGAresult}{变式}{4.8}
在 4.5 中去掉条件 (a)，令 $Z$ 为 $f$ 的非……点集。仍可得到
\[
\phi_x^i=0\qquad(i<n-\dim Z_s).
\]
\end{SGAresult}
对 $\dim\overline{\{x\}}$ 作\Term{descent}{下降}归纳。归纳假设使我们可以取 $m=n-\dim Z_s$ 应用 4.5 (1)，然后像前面一样完成证明。

\textbf{4.9.} 也可得到只利用 $f$ 的\Term{geometric-fibre}{几何纤维}信息的消没判据。不过在\Term{complete-intersection}{完全交}的情形，所得结果比 4.5 差一次。

\begin{SGAresult}{命题}{4.10}
在 4.5 的\Term{generic}{一般}假设下，假设：

\textup{(a)} $x$ 为其所在\Term{fibre}{纤维}中的孤立非\Term{smooth}{光滑}点；

\textup{(b)} 对 $X_{\bar\eta}$ 中每一个为 $x$ 的\Term{generalization}{一般化点}的点 $y$，若其在 $X_{\bar\eta}$ 中的\Term{closure}{闭包}\Term{dimension}{维数}为 $d(y)$，则
\[
\operatorname{prof\,\text{ét}}_y(X_{\bar\eta})\ge n-d(y),
\qquad\text{即}\quad
H^i_{\{y\}}(X_{\bar\eta})=0\quad(i\le n-d(y));
\]
\textup{(c)} $\operatorname{prof\,\text{ét}}_x(X_s)\ge n$。

则 $i<n-1$ 时，$\phi_x^i=0$。
\end{SGAresult}
保留 4.5 的记号；可以取 $m=\infty$ 应用 4.5 (1)。由\Term{local-scheme}{局部} Lefschetz \Term{theorem}{定理}（SGA 2 XIV），假设 (b) 给出
\[
\begin{aligned}
H^i(V')&\xrightarrow{\sim}H^i(V_s)&&(i<n-1),\\
H^i(V')&\lhook\joinrel\longrightarrow H^i(V_s)&&(i=n-1).
\end{aligned}
\]
再由 (c) 得到
\[
\widetilde H^i(\bar V)=\varinjlim\widetilde H^i(V')=0\qquad(i\le n-1),
\]
于是 4.10 得证。

\textbf{4.11.} 如同 4.8，也可给出 4.10 的一个不要求 (a) 的变式。

\SourcePage{26}
\SGAsection{5}{单值性对基本群的作用}
在本节（即 Grothendieck 于 1968 年 11 月 19 日所作的报告）中，我们将使用\Term{curve}{曲线}的\Term{semistable-reduction}{半稳定约化}\Term{theorem}{定理}。Artin 和 Winters [2] 给出了此\Term{theorem}{定理}的直接证明；它也可由\Term{abelian-variety}{阿贝尔簇}的相应\Term{theorem}{定理}（第 IX 讲第 3 节，使用 3.6）\Term{pushout}{推出} [3]。我们还将使用 Schlessinger 理论（Rim 的第 VI 讲）。

\textbf{5.1.} 设 $k$ 为\Term{algebraically-closed-field}{代数闭域}，$C$ 为 $k$ 上的\Term{projective}{射影}\Term{curve}{曲线}，$x_1,\ldots,x_n$ 为 $C$ 的\Term{finite}{有限}个\Term{closed-point}{闭点}。若满足下列条件，则称 $X=(C;x_1,\ldots,x_n)$ \Term{stable}{稳定}：

\textbf{(a)} $C$ \Term{reduced}{约化}，且其\Term{singular-point}{奇点}都为具有不同切线的\Term{double-point}{二重点}；

\textbf{(b)} 各 $x_i$ 互异，且都是 $C$ 的\Term{smooth}{光滑}点；

\textbf{(c)} 若 $C_1$ 为 $C$ 的一个\Term{irreducible-component}{不可约分支}，$E$ 为 $C_1$ 的\Term{normalization}{正规化} $C'_1$ 上位于 $C_1$ 的\Term{singular-point}{奇点}或某个 $x_i$ 上方的点的\Term{set}{集合}，则当 $C'_1$ 的\Term{genus}{亏格}为 $0$（分别为 $1$）时，$E$ 至少有 $3$（分别为 $1$）个元素。

条件 (c) 与下列每个条件\Term{equivalence}{等价}：

\textbf{(c')} $X$ 没有非平凡的\Term{infinitesimal-automorphism}{无穷小自同构}。

\textbf{(c'')} 若 $\omega$ 为\Term{dualizing}{对偶化}\Term{invertible-sheaf}{可逆层}，则\Term{invertible-sheaf}{可逆层} $\omega'=\omega(\sum x_i)$ 为\Term{ample}{丰沛}的。

若一个带\Term{section}{截面} $x_1,\ldots,x_n$ 的 $S$ \Term{scheme}{概形} $C$ 的\Term{geometric-fibre}{几何纤维}都是\Term{stable-curve}{稳定曲线}，则称其为 $S$ 上的\Term{stable-curve}{稳定曲线}。在 [3] 第 1、2 节中，只考虑了 $n=0$ 的情形。该文结果向\Term{generic}{一般}情形的推广是显然的：

\textbf{(1)} 若 $X/k$ \Term{stable}{稳定}，则 $n\ge2$ 时 $H^1(C,\omega'^{\otimes n})=0$，$n\ge3$ 时 $\omega'^{\otimes n}$ \Term{very-ample}{非常丰沛}（参见 [3] 1.2）。

\textbf{(2)} $X$ 的\Term{formal-moduli-scheme}{形式模概形} $M$ \Term{smooth}{光滑}，且 $C$ 保持\Term{singular}{奇异}的点位于 $M$ 的一个\Term{normal-crossings-divisor}{正规交叉除子}上。

\SourcePage{27}
\textbf{(3)} 若 $X$、$Y$ 在 $S$ 上\Term{stable}{稳定}，则 $\operatorname{Isom}_S(X,Y)$ 在 $S$ 上\Term{finite}{有限}且\Term{unramified}{非分歧}（这里的\Term{isomorphism}{同构}与标记点相容）。

\textbf{(4)} 若 $X/k$ \Term{stable}{稳定}，则 $\operatorname{Aut}_k(X)\to\operatorname{Aut}_k\operatorname{Pic}^0(C)$ 为\Term{injective}{单射}。

最后还有：
\begin{SGAresult}{命题}{5.1.1}
设 $X_\eta=(C_\eta;x_1,\ldots,x_n)$ 在\Term{trait}{离散赋值环谱} $S$ 的\Term{generic-point}{一般点} $\eta$ 上\Term{stable}{稳定}，且 $C_\eta$ \Term{smooth}{光滑}。存在 $\eta$ 的\Term{finite-separable-extension}{有限可分扩张} $\eta'$，以及 $S$ 在 $\eta'$ 中的\Term{normalization}{正规化} $S'$ 上的一条\Term{stable-curve}{稳定曲线} $X'$，使
\[
X'\otimes_{S'}\eta'=X_\eta\otimes_\eta\eta'.
\]
\end{SGAresult}
这由通常的\Term{semistable-reduction}{半稳定约化}\Term{theorem}{定理}得到：它允许选取 $\eta'$，使 $C_{\eta'}$ 在 $S'$ 上的\Term{minimal-model}{极小模型} $C''$ 的几何\Term{special-fibre}{特殊纤维}满足 (a)。然后在 $C''$ 的\Term{smooth}{光滑}点处作\Term{blowup}{爆破}，必要时反复进行，使 (b) 成立。再\Term{retraction}{收缩}\Term{special-fibre}{特殊纤维}中不包含任何 $x_i$、\Term{self-intersection}{自交数}为二的\Term{smooth}{光滑}有理\Term{branch}{分支}链，便得到 $C'$（参见 [3]，第 88 页 2.3 的证明）。

可像该文那样验证，只要 $C_\eta$ 的 Jacobian 簇有\Term{stable-reduction}{稳定约化}，$C'$ 就存在。此外，$C_\eta$ 的\Term{smoothness}{光滑性}假设是多余的。

\textbf{5.2.} 设 $X_\eta$ 为\Term{strictly-local}{严格局部}\Term{trait}{离散赋值环谱} $S$ 的\Term{fraction-field}{分式域}上的\Term{geometrically-connected}{几何连通}\Term{finite-type}{有限型}\Term{scheme}{概形}。对 $X_{\bar\eta}$ 的一个\Term{geometric-point}{几何点} $\xi$，有\Term{exact-sequence}{正合列}
\[
0\longrightarrow\pi_1(X_{\bar\eta},\xi)\longrightarrow\pi_1(X,\xi)
\longrightarrow\operatorname{Gal}(\bar\eta/\eta)\longrightarrow0,\tag{5.2.1}
\]
由此得到 $\operatorname{Gal}(\bar\eta/\eta)$ 到 $\pi_1(X_{\bar\eta})$ 的\Term{outer-automorphism}{外自同构}群的\Term{homomorphism}{同态}。转到 $\pi_1(X_{\bar\eta})$ 的最大的素于 $p$ 的\Term{quotient-group}{商群} $\pi_1^{(p)}(X_{\bar\eta})$，得到
\[
\operatorname{Gal}(\bar\eta/\eta)\longrightarrow
\operatorname{Aut\,ext}\bigl(\pi_1^{(p)}(X_{\bar\eta})\bigr).\tag{5.2.2}
\]
假设 $X_\eta$ 有一个\Term{rational-point}{有理点} $x$，令 $\bar x$ 为\Term{geometric-point}{几何点} $\bar\eta\to\eta\xrightarrow{x}X_\eta$。取 $\xi=\bar x$ 时，\Term{exact-sequence}{正合列} (5.2.2) \Term{canonical}{典范}\Term{split}{分裂}，因而有
\SourcePage{28}
\[
[x]:\operatorname{Gal}(\bar\eta/\eta)\longrightarrow
\operatorname{Aut}\bigl(\pi_1(X_{\bar\eta},\bar x)\bigr)
\longrightarrow\operatorname{Aut}\bigl(\pi_1^{(p)}(X_{\bar\eta},\bar x)\bigr).\tag{5.2.3}
\]
\begin{SGAresult}{\Term{theorem}{定理}}{5.3}
$P$ 在 5.2.2（或 5.2.3，两者\Term{equivalence}{等价}）下的像为\Term{finite-group}{有限群}。
\end{SGAresult}
\textbf{5.3.1. 归结到 $X_\eta$ \Term{geometrically-normal}{几何正规}且整的情形。}
设 $X'_\eta$ 为 $X_\eta$ 的\Term{normalization}{正规化}，$(X'_{i\eta})_i$ 为 $X'_\eta$ 的\Term{connected-component}{连通分支}族，$X''_\eta=X'_\eta\times_{X_\eta}X'_\eta$，$(X''_{j\eta})_j$ 为 $X''_\eta$ 的\Term{connected-component}{连通分支}。把 $k(\eta)$ 换成一个\Term{finite-extension}{有限扩张}后，可假设各 $X'_{i\eta}$ \Term{geometrically-normal}{几何正规}且整，并有\Term{rational-point}{有理点} $x_i$，各 $X''_{j\eta}$ \Term{geometrically-connected}{几何连通}，并有\Term{rational-point}{有理点} $y_j$。对 $k=1$ 或 $2$，以及 $p_k(y_j)\in X'_{i\eta}$，选取从 $\operatorname{pr}_k(\bar y_j)$ 到 $\bar x_i$ 的一个\Term{path}{道路}类 $\ell_{kij}$；这是一个 $\pi_1(X'_{i\bar\eta},\bar x_i)$ \Term{torsor}{挠子}的元素。令 $\bar\ell_{kij}$ 为相应的“素于 $p$ 的\Term{path}{道路}”（即模去 $\operatorname{Ker}(\pi\to\pi^{(p)})$ 后的同一\Term{path}{道路}）。可以选取被 $P$ 固定的 $\bar\ell_{kij}$：素于 $p$ 的\Term{path}{道路}\Term{set}{集合}是若干\Term{cardinal}{基数}素于 $p$ 的\Term{finite-set}{有限集合}的\Term{inverse-limit}{投射极限}，且 pro-$p$ 群 $P$ \Term{action}{作用}于这些\Term{finite-set}{有限集合}。

于是，若 $P$ 的某个\Term{finite-index}{有限指数}\Term{subgroup}{子群} $P'$ 在各 $\pi_1^{(p)}(X'_{i\bar\eta},\bar x_i)$ 上的\Term{action}{作用}平凡，则由 Van Kampen 公式（即 $X'_{\bar\eta}\to X_{\bar\eta}$ 对\Term{etale}{平展}\Term{covering}{覆盖}为\Term{effective-descent}{有效下降}\Term{morphism}{态射}）知，$P'$ 在 $\pi_1^{(p)}(X_{\bar\eta},\bar x)$ 上的\Term{action}{作用}也平凡。

\textbf{5.3.2.} 若 $X_{\bar\eta}$ \Term{normal}{正规}，$U$ 为其\Term{open-subset}{开集}，则 $\pi_1(U)$ \Term{surjective}{满射}到 $\pi_1(X_{\bar\eta})$；故可假设 $X_{\bar\eta}$ \Term{smooth}{光滑}且\Term{affine}{仿射}。用\Term{generic}{一般}\Term{hyperplane}{超平面}截取（这将 $S$ 换成某个 $\mathbf A_S^k$ 的\Term{special-fibre}{特殊纤维}的\Term{generic-point}{一般点}处的\Term{strict-localization}{严格局部化}），并应用 Bertini \Term{theorem}{定理}，便可归结到 $X_\eta$ \Term{smooth}{光滑}、\Term{geometrically-connected}{几何连通}且一维的情形。把 $k(\eta)$ 换成一个\Term{finite-extension}{有限扩张}后，可以假设 $X_\eta$ 是\Term{smooth}{光滑}\Term{projective}{射影}\Term{curve}{曲线} $\bar X_\eta$ 中\Term{finite}{有限}个\Term{rational-point}{有理点} $x_i$ 的补集（我们还可以增加这些点的个数）。应用\Term{semistable-reduction}{半稳定约化}\Term{theorem}{定理}（5.1.1），便可归结到下列情形：$S$ 上存在\Term{stable-curve}{稳定曲线} $X=(C;x_1,\ldots,x_n)$，使 $X_\eta=C_\eta-\{x_1,\ldots,x_n\}$。然后应用下述结果。

\SourcePage{29}
\begin{SGAresult}{\Term{theorem}{定理}}{5.4}
设 $S$ 为\Term{strictly-local}{严格局部}\Term{scheme}{概形}，$f:\bar X\to S$ 为\Term{proper}{固有}\Term{flat}{平坦}\Term{morphism}{态射}，其\Term{special-fibre}{特殊纤维}为一条\Term{curve}{曲线}，且非\Term{smooth}{光滑}点都为\Term{ordinary-double-point}{普通二重点}。设 $Y$ 为 $\bar X$ 的\Term{subscheme}{子概形}，是包含于\Term{smooth}{光滑}\Term{open-subset}{开集}中的\Term{finite}{有限}个互\Term{disjoint}{不相交}\Term{section}{截面}的并。设 $x$ 为 $X=\bar X-Y$ 的\Term{rational-point}{有理点}。设 $U$ 为 $S$ 中使 $\bar X$ 在其上\Term{smooth}{光滑}的\Term{open-subset}{开集}，$\xi$ 为 $U$ 的一个\Term{geometric-point}{几何点}。则 $\pi_1(U,\xi)$ 在 $\operatorname{Aut}(\pi_1^{(p)}(X_\xi,x_\xi))$ 中的像为\Term{abelian-group}{阿贝尔群}，且素于 $p$。
\end{SGAresult}
应注意，之所以能使 $\pi_1(U,\xi)$ \Term{action}{作用}于 $\pi_1^{(p)}(X_\xi,x_\xi)$，正是因为 $\pi_1^{(p)}$ 满足\Term{specialization}{特化}\Term{theorem}{定理}。

先归结到 $S$ \Term{complete}{完备}的情形，再归结到 $S$ 为带有各\Term{section}{截面}的\Term{special-fibre}{特殊纤维}的\Term{formal-moduli-scheme}{形式模概形}这一\Term{versal}{万有}情形。此时 $S$ 是 Cohen 环上的\Term{formal-power-series-ring}{形式幂级数环}的谱：$S=\operatorname{Spec}(W[[t_1,\ldots,t_n]])$，而 $U$ 是由方程 $t_1\cdots t_m=0$ 定义的\Term{normal-crossings-divisor}{正规交叉除子}的补集。在此\Term{versal}{万有}情形，由 Abhyankar \Term{lemma}{引理}，$\pi_1(U,\xi)$ 为\Term{abelian-group}{阿贝尔群}且素于 $p$，\Term{theorem}{定理}由此得证。

\textbf{注记 5.5。} 在任意\Term{relative-dimension}{相对维数}下，考虑一维的\Term{generic}{一般}平面\Term{section}{截面}并应用 Bertini \Term{theorem}{定理}，可证明类似于 5.4 的结果。

\SGAsection{6}{附录：稳定约化定理的算术证明}
\begin{center}P. Deligne\end{center}
设 $S$ 为\Term{trait}{离散赋值环谱}，$A_\eta$ 为 $k(\eta)$ 上\Term{dimension}{维数} $d\ne0$ 的\Term{abelian-variety}{阿贝尔簇}。对\Term{trait}{离散赋值环谱}间的\Term{local-scheme}{局部}\Term{morphism}{态射} $S'\to S$，以 $A_{\eta'}$ 表示由 $A_\eta$ 作\Term{scalar-extension}{标量扩张}得到的 $k(\eta')$ 上的\Term{abelian-variety}{阿贝尔簇}。设 $\mathcal A_{S'}$ 为 $A_{\eta'}$ 在 $S'$ 上的 \Term{neron-model}{Néron 模型}，$\mathcal A_{S's'}$ 为其\Term{special-fibre}{特殊纤维}，$\mathcal A^0_{S's'}$ 为 $\mathcal A_{S's'}$ 的\Term{identity-component}{单位分支}。\Term{stable-reduction}{稳定约化}\Term{theorem}{定理}陈述如下。

\SourcePage{30}
\begin{SGAresult}{\Term{theorem}{定理}}{6.1}
对适当的 $S'$，$A_{\eta'}$ 有\Term{stable-reduction}{稳定约化}，即 $\mathcal A^0_{S's'}$ 是\Term{abelian-variety}{阿贝尔簇}被\Term{torus}{环面}的\Term{group-extension}{扩张}。
\end{SGAresult}
由 6.1 很容易\Term{pushout}{推出}，可以选取 $S'$ 使 $k(\eta')$ 为 $k(\eta)$ 的\Term{finite-separable-extension}{有限可分扩张}。先假设 $S$ 的\Term{residue-field}{剩余域}在\Term{prime-field}{素域}上\Term{finite-generated}{有限生成}（若 $k(s)$ 是此种域的\Term{radicial}{根性}\Term{group-extension}{扩张}，同样的论证也适用）。

设 $\ell$ 为与剩余\Term{characteristic-exponent}{特征指数} $p$ 互素的\Term{prime-number}{素数}，$T_\ell(A)$ 为 \Term{tate-module}{Tate 模}，$V_\ell(A)=T_\ell(A)\otimes\mathbf Q_\ell$，$\rho:\operatorname{Gal}(\bar\eta/\eta)\to\operatorname{GL}(T_\ell(A))$ 为自然表示。先作\Term{scalar-extension}{标量扩张}，可归结到\Term{inertia-group}{惯性群} $I$ \Term{unipotent}{幺幂}\Term{action}{作用}的情形（1.2），于是其\Term{action}{作用}通过最大的 pro-$\ell$ \Term{quotient-group}{商群} $\mathbf Z_\ell(1)$（0.3）\Term{resolution}{分解}。设 $T$ 为 $\mathbf Z_\ell(1)$ 的一个\Term{generator}{生成元}在 $\operatorname{GL}(T_\ell(A))$ 中的像，$r$ 为使 $(T-1)^r\ne0$ 的最大整数 $\ge0$。有：
\begin{SGAresult}{\Term{lemma}{引理}}{6.2}
映射
\[
\sigma\in I,\quad v\in V_\ell(A)\longmapsto(\rho(\sigma)-1)^r(v)
\]
诱导下列\Term{morphism}{态射}及\Term{isomorphism}{同构}：
\[
\begin{tikzcd}[column sep=large]
V_\ell(A)_I\otimes\mathbf Q_\ell(r)\arrow[r,"M"]\arrow[d]&V_\ell(A)^I\\
\bigl(V_\ell(A)/\operatorname{Ker}(T-1)^r\bigr)\otimes\mathbf Q_\ell(r)\arrow[r,"\sim"]&\operatorname{Im}(T-1)^r\arrow[u]
\end{tikzcd}\tag{6.2.1}
\]
$V_\ell(A)_I$ 与 $V_\ell(A)^I$ 都是 $\operatorname{Gal}(\bar s/s)$ 模，且 $M$ 与 Galois \Term{action}{作用}相容。
\end{SGAresult}
\textbf{6.3.} 若 $\ell$ 进表示 $\tau:\operatorname{Gal}(\bar s/s)\to\operatorname{GL}(V)$ 满足下列条件，则称其\Term{weight}{权}为 $n$（$n\in\mathbf Z$）：

\textbf{($\star$)} 对 $k(s)$ 的某个适当\Term{model}{模型} $X$（$X$ 是\Term{prime-field}{素域}上的\Term{irreducible}{不可约}\Term{finite-type}{有限型}簇，且 $k(s)=k(X)$），$\tau$ 通过 $\pi_1(X,\bar s)$ \Term{resolution}{分解}，并且对 $X$ 的\Term{closed-point}{闭点} $x$，$\tau(F_x)$ 的\Term{eigenvalue}{特征值}为\Term{algebra}{代数}数，其所有复\Term{conjugacy}{共轭}的绝对值均为 $|k(x)|^{n/2}$（$F_x$ 表示几何 Frobenius $\phi_x^{-1}$）。

\SourcePage{31}
若 $V$ 的\Term{weight}{权}为 $n$，$W$ 的\Term{weight}{权}为 $m$，且 $n\ne m$，则 $\operatorname{Hom}_{\mathrm{Gal}}(V,W)=0$。

\textbf{6.4.} 由 \Term{neron-model}{Néron 模型}的\Term{universal-property}{泛性质}，$V_\ell(A)^I=V_\ell(\mathcal A^0_{S,s})$。设 $a$ 为 $\mathcal A^0_{S,s}$ 的最大\Term{abelian-quotient}{阿贝尔商}的\Term{dimension}{维数}，$m$ 为 $\mathcal A^0_{S,s}$ 的\Term{multiplication}{乘法}部分的\Term{dimension}{维数}。由 Weil 的结果，$V_\ell(A)^I$ 因而是一个\Term{dimension}{维数} $2a$、\Term{weight}{权}为 $-1$ 的表示被一个\Term{dimension}{维数} $m$、\Term{weight}{权}为 $-2$ 的表示的\Term{group-extension}{扩张}。

\textbf{6.5.} 设 $A^*$ 为 $A$ 的\Term{duality}{对偶}\Term{abelian-variety}{阿贝尔簇}。回顾 $V_\ell(A^*)$ 与 $V_\ell(A)$（因而 $V_\ell(A^*)^I$ 与 $V_\ell(A)_I$）之间有取值于 $\mathbf Q_\ell(1)$ 的\Term{duality}{对偶}。若 $a^*$、$m^*$ 对 $A^*$ 的定义分别与 $a$、$m$ 对 $A$ 的定义相同，则由 6.4，$V_\ell(A)_I$ 是一个\Term{dimension}{维数} $m^*$、\Term{weight}{权}为 $0$ 的表示被一个\Term{dimension}{维数} $2a^*$、\Term{weight}{权}为 $-1$ 的表示的\Term{group-extension}{扩张}。

\textbf{6.6.} Galois 模 $\mathbf Q_\ell(r)$ 的\Term{weight}{权}为 $-2r$。由于 $\operatorname{Im}(T-1)^r\ne0$，由 6.4、6.5 及 6.2 中的\Term{isomorphism}{同构}\Term{pushout}{推出} $r\le1$。若 $r=0$，则有
\[
\begin{cases}a+m\le d,\\2a+m=2d,\end{cases}
\]
故 $m=0$、$a=d$，$A_\eta$ 有\Term{good-reduction}{良约化}。再设 $r=1$。由 6.2 得到 $\operatorname{Im}(T-1)$ 的\Term{weight}{权}为 $-2$，而 $V_\ell(A)/V_\ell(A)^I$ 的\Term{weight}{权}为 $0$；这两个\Term{space}{空间}有相同的\Term{dimension}{维数} $m_1\le m$。

于是有
\[
\dim V_\ell(A)^I=2d-m_1=2a+m,
\]
\[
2\dim\mathcal A_{S,s}\ge2(a+m)\ge2a+m+m_1=2d=2\dim\mathcal A_{S,s},
\]
故 $A_\eta$ 有\Term{stable-reduction}{稳定约化}（$\dim\mathcal A_{S,s}=a+m$）。同时还得到 $m_1=m$，即
\[
\operatorname{Im}(T-1)\subset V_\ell(A)^I\simeq V_\ell(\mathcal A^0_{S,s})
\]
正是 $\mathcal A^0_{S,s}$ 的\Term{multiplication}{乘法}部分的 $V_\ell$。

\textbf{6.7.} 为处理\Term{generic}{一般}情形，我们将使用 0.5（参见 1.3）。可归结到假设：
\SourcePage{32}
\textbf{(a)} (0.5.2) 或 (0.5.3) 的条件成立；

\textbf{(b)} $\operatorname{Gal}(\bar\eta/\eta)$ 在 $A_\ell$ 上的\Term{action}{作用}平凡；

\textbf{(c)} $I$ 在 $T_\ell(A)$ 上的\Term{action}{作用}\Term{unipotent}{幺幂}。

在这些假设下，我们将证明 $A_\eta$ 有\Term{stable-reduction}{稳定约化}。由假设 (a)，可应用 (0.5.1)，像 1.3 那样得到 $S=\varprojlim S_i$；沿用该处的记号。设 $s_i$ 为 $S_i$ 的\Term{closed-point}{闭点}。

当 $i$ 充分大时，\Term{neron-model}{Néron 模型} $\mathcal A_S$ 的\Term{identity-component}{单位分支} $\mathcal A^0$ 来自 $S_i$ 上的一个\Term{fibre}{纤维}\Term{connected}{连通}的\Term{smooth}{光滑}\Term{group-scheme}{群概形} $\mathcal A_i$，且 $(\mathcal A_i)_\ell$ 在 $S_i-D_i$ 上为\Term{constant}{常值}。群 $\pi_1(S_i-D_i,\bar\eta)$ \Term{action}{作用}于 $T_\ell(A)$，并在 $T_\ell(A)/\ell T_\ell(A)\simeq A_\ell$ 上\Term{trivial-action}{平凡作用}。如同 1.3，得到 $P_i$ 的\Term{action}{作用}平凡。于是由 (1.3.1)：
\[
\begin{gathered}
T_\ell(A)_I=T_\ell(A)_{I_i},\\
T_\ell(\mathcal A_{i,s_i})=T_\ell(\mathcal A_s^0)
=T_\ell(A)^I=T_\ell(A^{I_i}),
\end{gathered}
\]
而 (6.2.1) 是 $\operatorname{Gal}(\bar s_i/s_i)$ 模的图，因而可对其应用 6.4 至 6.6 的论证。

\SourcePage{33}
\SGAsection{}{参考文献}
\begin{description}
\item[{[1]}] M. Artin，\emph{Algebraic approximation of structures over complete local rings}（\Term{complete}{完备}\Term{local-ring}{局部环}上结构的\Term{algebra}{代数}逼近），Publ. Math. I.H.E.S. \textbf{36}，1969，第 23—58 页。
\item[{[2]}] M. Artin、G. Winters，\emph{Degenerate fibres and reduction of curves}（\Term{degeneration}{退化}\Term{fibre}{纤维}与\Term{curve}{曲线}的\Term{reduced}{约化}）。
\item[{[3]}] P. Deligne、D. Mumford，\emph{The irreductibility of the space of curves of a given genus}（给定\Term{genus}{亏格}的\Term{curve}{曲线}\Term{space}{空间}的\Term{irreducible}{不可约}性），Publ. Math. I.H.E.S. \textbf{36}，1969，第 75—110 页。
\item[{[4]}] J.-P. Serre，\emph{Corps locaux}（\Term{local-field}{局部域}），Publ. Math. univ. Nancago–Hermann，1962。
\item[{[5]}] J.-P. Serre、J. Tate，\emph{Good reduction of abelian varieties}（\Term{abelian-variety}{阿贝尔簇}的\Term{good-reduction}{良约化}），Ann. of Math. \textbf{88}，第 3 期，1968，第 492—517 页。
\end{description}
